\documentclass[10pt,twoside]{article}
\newcommand{\BedrockArchitectureName}{Bedrock}
\newcommand{\BedrockDocumentTitle}{Bedrock C ABI}
\newcommand{\BedrockRunningTitle}{C ABI}
\newcommand{\BedrockFooterTitle}{BEDROCK ARCHITECTURE}
\usepackage{bedrock-reference}
\begin{document}
\BedrockMakeTitlePage{Bedrock C ABI}{C Language Binding and Calling Convention}{Draft}
\BedrockFrontMatter

\clearpage
\section{Scope}\label{page:scope}
This document defines the Bedrock C data model, register convention, calling
convention, aggregate and variadic rules, compiler-runtime helper ABI, and C
memory-model lowering. The Bedrock ELF ABI and ISA supply the referenced
binary-format and architectural objects. Target-specific C extensions own
their source syntax and semantics.

\subsection{Contract Position}
\BedrockContractDiagram{language}

\subsection{Profile}
\BedrockField{Profile:}{\texttt{bedrock-c}}

\clearpage
\section{Terminology}\label{page:terminology}
These terms are specific to the C binding. ELF and address-space terms keep the meanings defined by the Bedrock ELF ABI and the architecture manual.
\begin{description}[style=nextline,leftmargin=1.45in,labelwidth=1.35in,itemsep=3pt,topsep=3pt]
\item[object pointer] A stable 64-bit pre-segment coordinate that designates an object or one-past an object according to the C abstract machine.
\item[function pointer] A stable 64-bit pre-segment coordinate naming a C entry point.
\item[caller-saved register] A register whose live incoming value the caller preserves across a call.
\item[callee-saved register] A register whose incoming value must be preserved by the called function if the function uses it.
\item[sret buffer] Caller-allocated storage used to receive a return value that is too large or unsuitable for the normal return registers.
\item[addressable by-value argument] A by-value aggregate argument that is passed through caller-owned stack storage because the callee may need an address for the object.
\item[scalable value] A target vector or predicate value whose complete object
size is determined by the implementation-selected MAX\_VLEN. Source-language
spellings belong to the applicable compiler extension; this calling convention
defines their binary interface.
\end{description}

\clearpage
\section{Data Model}\label{page:data-model}
\BedrockField{Model:}{LP64}
\BedrockField{Byte Order:}{little endian}
\BedrockField{Internal SP:}{8-byte aligned}
\BedrockField{Function Entry:}{2-byte aligned}
\BedrockField{Before CALL:}{SP mod 16 = 8}
\BedrockField{Aggregate Maximum Alignment:}{16 bytes}
\subsection{Scalar Types}
\BedrockTableCaption{Reference C Data Model}
\begingroup\footnotesize
\setlength{\tabcolsep}{2pt}
\begin{longtable}{@{}p{2.1in}p{1.0in}p{1.0in}@{}}
\toprule
\textbf{Type} & \textbf{Bits} & \textbf{Align}\\
\midrule
\endhead
\texttt{\_\allowbreak{}Bool} & 8 & 1\\
\texttt{char} & 8 & 1\\
\texttt{signed char} & 8 & 1\\
\texttt{unsigned char} & 8 & 1\\
\texttt{short} & 16 & 2\\
\texttt{int} & 32 & 4\\
\texttt{long} & 64 & 8\\
\texttt{long long} & 64 & 8\\
\texttt{\_\allowbreak{}\_\allowbreak{}int128} & 128 & 16\\
\texttt{unsigned \_\allowbreak{}\_\allowbreak{}int128} & 128 & 16\\
\texttt{ordinary pointer} & 64 & 8\\
\texttt{size\_\allowbreak{}t} & 64 & 8\\
\texttt{ptrdiff\_\allowbreak{}t} & 64 & 8\\
\texttt{wchar\_\allowbreak{}t} & 32 & 4\\
\texttt{char16\_\allowbreak{}t} & 16 & 2\\
\texttt{char32\_\allowbreak{}t} & 32 & 4\\
\texttt{float} & 32 & 4\\
\texttt{double} & 64 & 8\\
\texttt{long double} & 64 & 8\\
\texttt{float \_Complex} & 64 & 4\\
\texttt{double \_Complex} & 128 & 8\\
\texttt{long double \_Complex} & 128 & 8\\
\bottomrule
\end{longtable}
\endgroup
\subsection{Scalar Type Semantics}
The widths and natural alignments of scalar types are given in
Table~3-1. The following rules define the type identities and representations
that supplement those two values.

\paragraph{Integer types and aliases.}
Plain \texttt{char} is signed. \texttt{size\_t} is an alias for
\texttt{unsigned long}, and \texttt{ptrdiff\_t} is an alias for
\texttt{long}. \texttt{wchar\_t} is a signed 32-bit integer;
\texttt{char16\_t} and \texttt{char32\_t} are unsigned 16-bit and 32-bit
integers, respectively. In the absence of a separate enum extension, an enum
has the representation, size, and alignment of \texttt{int}.

A stored \texttt{\_Bool} has the integer value zero or one. False is encoded
with every object-representation bit clear; true is encoded with bit 0 set and
every other bit clear.

\paragraph{Ordinary pointers.}
An ordinary object or function pointer is 64 bits wide and 8-byte aligned. Its
value is a pre-segment coordinate. Object and function pointers remain distinct C type categories
even though their baseline representations have the same size and alignment.
The null pointer has numeric address zero and an all-zero
64-bit object representation. Every ordinary-pointer result with numeric
address zero is the null pointer.

\paragraph{128-bit integers.}
The types \texttt{\_\_int128} and \texttt{unsigned \_\_int128} are 16-byte,
16-byte-aligned integers. Their little-endian object representation places the
low 64-bit word at the lower address. When the calling convention assigns one
to a GENERAL-PAIR, the low word occupies the lower-numbered even register and
the high word occupies the following odd register. A returned 128-bit integer
therefore uses \texttt{R0} for the low word and \texttt{R1} for the high word.
Its stack slot is 16-byte aligned.

An ordinary 128-bit load or store may use two 64-bit memory operations. The
baseline C atomic object set contains the naturally aligned 1-, 2-, 4-, and
8-byte integer and pointer types.

\paragraph{Floating-point profile.}
The \texttt{bedrock-c} profile requires the base floating-point extension.
Floating-point helpers may replace
unavailable operations only when the \texttt{F0}--\texttt{F15} register calling
convention remains available.

\paragraph{Long double.}
\texttt{long double} is a distinct C type with the IEEE-754 binary64 object
representation, 8-byte size, and 8-byte alignment. It uses the floating-point
argument class and returns in \texttt{F0}. It therefore has the same call
representation as \texttt{double}. In a variadic call it occupies one
16-byte slot containing its 8-byte binary64 representation.

\paragraph{Complex types.}
A complex object contains its real component first and its imaginary component
second, each using the corresponding real type's representation. Therefore
\texttt{float \_Complex} has size 8 and alignment 4, while
\texttt{double \_Complex} and \texttt{long double \_Complex} have size 16 and
alignment 8.

\paragraph{Scalable vector and predicate types.}
When the scalable-vector extension is present, let $B$ be MAX\_VLEN/8 bytes. A
vector object contains exactly $B$ bytes and has 16-byte alignment. A
predicate object contains exactly $B/8$ packed-image bytes. A private vector
spill slot has $B$ bytes, 16-byte alignment, and $B$-byte stride. A private
predicate spill slot has 16-byte alignment and
$\operatorname{align\_up}(B/8,16)$-byte stride. Predicate spill tail padding
has unspecified contents. MAX\_VLEN is stable until reset, so every object and
frame in one execution context uses one size. VLEN may change while those
objects and frames exist; it controls active execution width without changing
their complete representations.

\subsection{Aggregate Layout}
Aggregate layout applies the scalar and pointer sizes and alignments defined
above as follows:
\begin{itemize}
\item Struct fields are laid out in declaration order with natural alignment up to the aggregate maximum.
\item A union's alignment is the greatest alignment required by any member. Its size is the greatest ABI size of any
member, rounded up to a multiple of the union alignment.
\item Arrays store elements contiguously with each element using the ABI size of its type.
\item Aggregate object size is rounded up to the aggregate object's alignment.
\end{itemize}

\paragraph{Bit-fields.}
A bit-field allocation unit has the size and alignment of its declared base
type. Within the little-endian unit, allocation begins at the low bit. A field
remains within one allocation-unit boundary, and only consecutive fields whose
base types are identical after typedef expansion may share a unit. A base-type
change or insufficient remaining space closes the current unit and begins a new
naturally aligned unit. A field width is at most the representation width of
its base type.

An externally visible bit-field base type shall be \texttt{\_Bool},
\texttt{signed int}, \texttt{unsigned int}, or a typedef of one of those types.
A plain \texttt{int} bit-field has the value range and representation of
\texttt{signed int}. The baseline ABI does not define an external calling or
object-layout contract for bit-fields with any other base type.

An unnamed nonzero-width field consumes space normally. A zero-width field
closes the current unit and advances layout to the next boundary for its base
type while contributing zero storage bytes. Ordinary aggregate alignment,
tail-padding, and size-rounding rules still apply after bit-field allocation.

\paragraph{Flexible arrays and extension layout facilities.}
A flexible array member contributes zero element-storage bytes to \texttt{sizeof} its
containing structure. Its offset is aligned for the element type, and that
element alignment participates in the structure's alignment calculation.

For an external object, a supported source-language packing annotation is part
of the declared object layout. A record-level \texttt{packed} annotation caps
each member's placement alignment and the record alignment at one byte. A
\texttt{\#pragma pack} limit caps each affected member's placement alignment
and the record alignment at the declared limit. A field-level packing
annotation caps that field's placement alignment at one byte. Fields remain in
declaration order, unions still overlay their members, and the aggregate size
is rounded to its resulting alignment. An array of such a type stores elements
contiguously with stride \texttt{sizeof} its element type. These rules apply
identically to an external definition and every compatible declaration.

Packing changes placement and permitted access alignment; it does not change a
member's value representation. Loads and stores through a packed member use
the member's effective alignment. Empty structures and zero-length arrays have
no external calling or object-layout contract in this ABI.

\section{Register Convention}\label{page:register-convention}
This section defines which architectural state survives an ordinary C call.
Argument and result locations are defined by the assignment procedure in
Section~5; that procedure is their sole definition.

\subsection{Register Preservation}
The caller shall preserve any live value held in a volatile register before a
call. A callee that modifies a nonvolatile register shall restore its exact
incoming value before returning. Volatile registers carry explicitly assigned
call results and caller-managed live values across the call.

\BedrockTableCaption{C Call Preservation Sets}
\begin{BedrockLongTable}{@{}p{1.35in}p{2.0in}p{2.05in}@{}}
\toprule
\textbf{State} & \textbf{Volatile across a call} & \textbf{Preserved by the callee}\\
\midrule
\endhead
General registers & \texttt{R0}--\texttt{R7} & \texttt{R8}--\texttt{R15}\\
Floating-point registers & \texttt{F0}--\texttt{F7} & \texttt{F8}--\texttt{F15}\\
Vector registers & \texttt{V0}--\texttt{V15} & \texttt{V16}--\texttt{V31}\\
Predicate registers & \texttt{P0}--\texttt{P7} & \texttt{P8}--\texttt{P15}\\
Vector status & none & \texttt{VSTATUS}\\
Condition state & \texttt{FLAGS} & none\\
Segment state & \texttt{GS1}--\texttt{GS5} & \texttt{CS}, \texttt{DS}, \texttt{SS}, \texttt{GS0}\\
\bottomrule
\end{BedrockLongTable}

A callee may preserve floating-point registers with \texttt{FPUSHP} and
\texttt{FPOPP} or with any sequence having the same observable effect. Each
instruction operates on one canonical pair. The nonvolatile pairs are index 4
for \texttt{F8:F9}, index 5 for \texttt{F10:F11}, index 6 for
\texttt{F12:F13}, and index 7 for \texttt{F14:F15}. When several pairs are
saved, a normal prologue pushes them in increasing pair-index order and the
matching epilogue pops them in decreasing pair-index order. A callee need only
save the pair containing each nonvolatile register it modifies.

A callee that modifies a nonvolatile vector or predicate register shall
restore that register's complete scalable image before returning. Saving only
the currently active lanes, only significant predicate bits for one element
size, or a fixed-width prefix represents a partial save; preservation requires
the complete scalable image.

A callee shall restore the exact incoming \texttt{VSTATUS} image before
returning. This preserves the caller's selected active vector length.

\subsection{Stack and Control State}
\texttt{SP} is the dedicated architectural stack pointer. A
callee may move it while executing but shall restore the entry value before
returning; Section~5 defines the additional alignment and frame rules.
\texttt{PC} is the architectural program counter and is visible to C code
only as a control-flow target. The call and return instructions define its
transition across a call. Condition codes in
\texttt{FLAGS} are volatile and may be clobbered by every callee.

\subsection{Segment Context}
An ordinary C call preserves the exact incoming images of \texttt{CS},
\texttt{DS}, \texttt{SS}, and \texttt{GS0}. The first three establish the
code, ordinary-data, and stack contexts. \texttt{GS0} establishes the TLS
context and may use translated-window mode. \texttt{GS1} through \texttt{GS5}
are caller-saved scratch segment registers; callers preserve any live incoming
values they hold.

\subsection{C Address-Space Context}
Ordinary object and function pointers are stable 64-bit pre-segment
coordinates. Every C ABI execution path maintains the following invariants:

\begin{itemize}
\item While a pointer is live, calls, returns, register or memory preservation,
context transitions, and ABI boundaries preserve its 64-bit coordinate.
\item Every ABI-permitted segment path that dereferences a valid object
pointer selects the same byte of the same C object for a given pointer value.
This includes a stack object's address after it escapes to an ordinary data
path, and a materialized TLS object's address.
\item Pointer equality, subtraction, and in-object arithmetic have the same C
result on every ABI-permitted path. A one-past pointer participates in those
operations and is an arithmetic boundary sentinel. Storage accesses use
coordinates within the object.
\item Function pointers, call entries, return continuations, ELF symbols and
relocations, dynamic linkage, stack pointers, runtime-provided pointers, TLS
materialization, signal and nonlocal-jump contexts, unwind records, and
debugger-visible coordinates use this pre-segment coordinate system.
\item Segment and page-table context transitions preserve these invariants as
one architectural transition from the program's perspective.
\item Every C object and its one-past coordinate must be representable as a
nonzero 64-bit address, including its one-past coordinate. The greatest
representable one-past coordinate is $2^{64}-1$.
\end{itemize}

The 64-bit all-zero null-pointer representation designates the null pointer.
Every C object and function occupies a nonzero address.

\subsection{Status Registers}
Target-specific intrinsics and runtime code provide C access to
\texttt{STATUS}. Low-level access to
\texttt{FSTATUS} and \texttt{FFLAGS} uses
\texttt{<bedrockfpuintrin.h>}. \texttt{FSTATUS} is callee-saved: a function
returns with the incoming rounding mode, floating-point exception-condition enables, and other control fields.
A documented floating-point-environment interface instead specifies its
outgoing environment explicitly.
\texttt{FFLAGS} is cumulative floating-point status and is caller-clobbered;
floating-point operations in a callee may set accrued floating-point exception flags. An
interface that requires a saved environment uses its floating-point environment
runtime contract explicitly.

\section{Calling Convention}\label{page:calling-convention}
This section is normative. A caller and callee conform when they independently
apply the classification and assignment procedure below and therefore agree on
every argument location, return location, and stack boundary. Register lists
are summaries; the assignment procedure remains authoritative.

\subsection{Stack and Calls}
The stack grows downward. An architectural \texttt{CALL} pushes an
8-byte return address at \texttt{[SP]}, and
callers align \texttt{SP} before the call so callee entry observes
16-byte alignment. Between ABI call boundaries, a function
may adjust \texttt{SP} while keeping it aligned to at least
8 bytes. Immediately before \texttt{CALL}, the stack
pointer satisfies SP modulo 16 equals 8 before return address push. After \texttt{CALL} pushes the
return address, it satisfies SP modulo 16 equals 0 at callee entry.

Each object observes its declared type alignment. In particular, relaxing
the internal stack-pointer invariant preserves the 16-byte alignment of
\texttt{\_\_int128} objects or other objects whose declared alignment is 16
bytes. Both the red-zone size and fixed outgoing-argument-area size are zero. A
callee that dynamically realigns the stack restores the incoming stack
pointer before returning.

The first stack argument begins at \texttt{[SP+16]} at callee entry. The eight
bytes at \texttt{[SP+8]} are alignment padding. Argument values begin at
\texttt{[SP+16]}.
Each subsequent stack argument occupies the next 16-byte slot.

\begin{BedrockListedStackFrameDiagram}{Near-call stack frame at callee entry}
\BedrockStackFrameSPSlot{+0}{saved return \texttt{PC}}
\BedrockStackFrameSlot{+8}{8-byte alignment padding}
\BedrockStackFramePayload{+16}{first stack argument, if present}
\end{BedrockListedStackFrameDiagram}

\subsection{Frame Pointer}
A function that maintains a frame pointer uses callee-saved \texttt{R15}.
Other functions retain \texttt{R15} as an ordinary callee-saved register.

\subsection{C Unwind Frame}
The Bedrock ELF ABI owns the language-independent DWARF register-number
assignment. C unwind descriptions use that assignment with the following
calling-convention frame rule.
At call entry, the canonical frame address is entry \texttt{SP+8}; saved
\texttt{PC} is at CFA minus 8 and \texttt{CS} has a same-value rule.

\Needspace{2.4in}
\subsection{Argument Classification}
Classification occurs after the C language has determined the effective type
of the argument. A named argument to a prototyped function is placed in one of
seven ABI classes:

\begin{description}[style=nextline,leftmargin=1.35in,labelwidth=1.25in,itemsep=3pt,topsep=3pt]
\item[GENERAL] Integer types up to 64 bits wide and ordinary object or
function pointers. One general-register location is required.
\item[GENERAL-PAIR] Signed and unsigned \texttt{\_\_int128}. Two consecutive
general registers beginning at an even-numbered register are required.
\item[FLOAT] \texttt{float}, \texttt{double}, and \texttt{long double}. One
floating-point register is required.
\item[FLOAT-PAIR] \texttt{float \_Complex}, \texttt{double \_Complex}, and
\texttt{long double \_Complex}. Two consecutive floating-point registers are
required with unit register alignment.
\item[VECTOR] A direct scalable vector value. One vector register is required.
\item[PREDICATE] A direct scalable predicate value. One predicate register is
required.
\item[INDIRECT] Every structure or union argument, regardless of size or
member composition. The caller creates a 16-byte-aligned by-value copy and
passes the address of that copy as one GENERAL value.
\end{description}

A compound scalable value is INDIRECT. Exhaustion of a scalable value's
direct register class also converts it to INDIRECT: the caller creates an
object with the layout above and assigns its pointer through the GENERAL
procedure. Each scalable value travels wholly in its register or through an
indirect copy.

Signed GENERAL values narrower than 64 bits are sign-extended when assigned to
a register. Unsigned values and \texttt{\_Bool} are zero-extended. Plain
\texttt{char} follows its signed ABI type. A stack slot contains the effective
type's object representation at its low address; bytes after that
representation are unspecified.

\subsection{Assignment Procedure}
The caller shall apply the following procedure. The callee may rely on the
result solely from assigned registers and value bytes.

\begin{enumerate}[leftmargin=1.35em,itemsep=3pt,topsep=3pt]
\item If the result uses an sret buffer, assign its address to \texttt{R0} and
initialize the general-register cursor to 1. Otherwise initialize it to 0.
Initialize the floating-point cursor to 0, the vector cursor to 0, and the
predicate cursor to 0.
\item Visit arguments in source order. GENERAL and INDIRECT arguments consume
the register named by the current general cursor and then advance that cursor.
FLOAT arguments analogously consume \texttt{F0} through \texttt{F7} using the
independent floating-point cursor. VECTOR arguments consume \texttt{V0}
through \texttt{V7} using the independent vector cursor, and PREDICATE
arguments consume \texttt{P0} through \texttt{P3} using the independent
predicate cursor.
\item For a FLOAT-PAIR argument, use the current floating-point register for
the real component and the next register for the imaginary component, then
advance the cursor by two with unit register alignment. If fewer than two
registers remain, place the complete complex value in one stack
slot and mark the floating-point class exhausted.
\item For a GENERAL-PAIR argument, round the general cursor upward to an even
register number. If that register and its successor both exist, place the low
64 bits in the even register, the high 64 bits in the odd register, and advance
the cursor past the pair. A skipped odd register remains unused.
\item If the remaining registers of a class have insufficient capacity for an
argument, place the complete argument in one stack slot and mark
that register class exhausted. Later arguments of that class also use stack
slots, while register holes remain unassigned. Each argument occupies one
complete register location or one complete stack slot.
For VECTOR or PREDICATE, instead convert the complete argument to INDIRECT as
specified above and assign that pointer through GENERAL; exhausting a scalable
register class leaves GENERAL independently available.
\item Unnamed variadic arguments and every argument of an unprototyped call are
forced to stack slots and leave register cursors unchanged. Apply the default
argument promotions before assigning those slots.
Every scalable argument in either case is first converted to INDIRECT. Its
GENERAL pointer is then subject to this forced-stack rule, so the variadic slot
contains the pointer exclusively.
\item Assign stack slots to stack-bound arguments in source order beginning at
\texttt{[SP+16]}. Consequently, a caller that constructs the outgoing area by
stores or pushes does so from the rightmost stack-bound argument toward the
leftmost. The reserved register-home-slot count is zero.
\end{enumerate}

The general, floating-point, vector, and predicate resources exhaust
independently.

\subsection{C Return Values}
Scalar results use their natural register class. A \texttt{\_\_int128} result
places its low word in \texttt{R0} and high word in \texttt{R1}. A structure
or union of at most 16 bytes is returned as raw object bytes: increasing object
addresses map first to \texttt{R0} and then to \texttt{R1}. Bytes beyond the
object size and padding bytes with C-unspecified values are
unspecified.

A complex result places its real component in \texttt{F0} and its imaginary
component in \texttt{F1}.

A single direct vector result uses \texttt{V0}; a single direct predicate
result uses \texttt{P0}. A compound scalable result or any result containing
multiple scalable values uses the ordinary indirect-result procedure, even if
its fixed fields would otherwise fit in \texttt{R0:R1}.

For a larger aggregate, the caller allocates a 16-byte-aligned result buffer
and passes its address in \texttt{R0} before ordinary argument assignment. The
callee stores the result in that buffer and returns the same address in
\texttt{R0}.

(:c-return-rules:base.value_classes.GENERAL_SCALAR,base.value_classes.GENERAL_PAIR,base.value_classes.FLOAT_SCALAR,base.value_classes.FLOAT_PAIR,base.value_classes.VECTOR_VALUE,base.value_classes.PREDICATE_VALUE,base.value_classes.AGGREGATE,base.value_classes.COMPOUND_SCALABLE:)

\subsection{Aggregate Passing}
For every by-value aggregate argument, the caller allocates distinct storage,
aligns its start to 16 bytes, initializes it with the argument value, and keeps
it alive until the call returns. The GENERAL procedure assigns the address
exclusively. If GENERAL registers are exhausted, the
stack argument slot contains that address exclusively.

The copy gives the callee an addressable object with the declared aggregate
type. Modifications to the private copy remain isolated from the caller's
source object. Baseline aggregates containing bit-fields and mixed
integer/floating aggregates use the same INDIRECT rule; their member
composition preserves the INDIRECT classification. Passing or returning an
aggregate with packed or unaligned members by value is outside the external
calling contract. This restriction does not apply to an external packed object
or to passing that object's address.

\subsection{Variadic Calls}
Fixed named arguments use the ordinary assignment procedure. Every unnamed
argument is first subject to the C default argument promotions and is then
placed in a 16-byte stack slot. In particular, \texttt{float} becomes
\texttt{double}, and \texttt{\_Bool}, \texttt{char}, \texttt{short}, and their
signed or unsigned variants promote to \texttt{int} when \texttt{int} can
represent all values of the source type. \texttt{long double} remains the
Bedrock binary64 \texttt{long double} type. An aggregate slot contains the
address of the caller-owned by-value copy described above.

The baseline \texttt{va\_list} representation is one 64-bit pointer to the next
unnamed argument slot. \texttt{va\_start} initializes it to the first unnamed
slot after any stack-assigned named arguments. \texttt{va\_arg} reads the
promoted representation from the low bytes of the current slot, or follows the
copy address for an aggregate, and then advances the pointer by 16 bytes.
\texttt{va\_copy} copies the pointer value and \texttt{va\_end} has a
zero-instruction machine effect.

An unprototyped call applies the same promotions and places every argument on
the stack. As required by C, the callee type must be compatible with the
promoted types.

\section{Freestanding C Binding}\label{page:freestanding-c-binding}
This section defines the compiler-runtime and ELF linkage obligations of the
freestanding C profile. It covers compiler-generated helper calls, C symbols,
and ELF linkage.

\subsection{Compiler Runtime Helpers}
The compiler lowers addition, subtraction, negation, bitwise operations,
comparisons, truncation, and extension of 128-bit integers inline. It may call
one of the helpers below for 128-bit multiplication,
division, remainder, and shifts; conversion between a 128-bit integer and a
supported floating-point type; or complex multiplication and division.

Every compiler-runtime helper is an ordinary \texttt{bedrock-c} function. Its
arguments and results therefore follow Section~5, including the
\texttt{R1:R0} return convention for \texttt{\_\_int128}.

The tables below are the complete baseline helper set. A compiler emits
implicit references from this set or from a target extension that defines the
symbol and its ABI. A runtime may export additional symbols as explicit
interfaces.

In the signature column, \texttt{i128} means \texttt{\_\_int128},
\texttt{u128} means \texttt{unsigned \_\_int128}, \texttt{c32} means
\texttt{float \_Complex}, and \texttt{c64} means
\texttt{double \_Complex}. The signatures use ordinary C parameter and result
classification; a shift count has type \texttt{int}.

\Needspace{2.7in}
\BedrockTableCaption{\texttt{\_\_int128} Arithmetic and Shift Helpers}
\begingroup\footnotesize
\setlength{\tabcolsep}{2pt}
\begin{longtable}{@{}p{1.55in}p{1.55in}p{2.45in}@{}}
\toprule
\textbf{Operation} & \textbf{Helper} & \textbf{Signature}\\
\midrule
\endhead
Multiply & \texttt{\_\allowbreak{}\_\allowbreak{}multi3} & \texttt{i128(i128, i128)}\\
Signed divide & \texttt{\_\allowbreak{}\_\allowbreak{}divti3} & \texttt{i128(i128, i128)}\\
Unsigned divide & \texttt{\_\allowbreak{}\_\allowbreak{}udivti3} & \texttt{u128(u128, u128)}\\
Signed remainder & \texttt{\_\allowbreak{}\_\allowbreak{}modti3} & \texttt{i128(i128, i128)}\\
Unsigned remainder & \texttt{\_\allowbreak{}\_\allowbreak{}umodti3} & \texttt{u128(u128, u128)}\\
Shift left & \texttt{\_\allowbreak{}\_\allowbreak{}ashlti3} & \texttt{i128(i128, int)}\\
Arithmetic shift right & \texttt{\_\allowbreak{}\_\allowbreak{}ashrti3} & \texttt{i128(i128, int)}\\
Logical shift right & \texttt{\_\allowbreak{}\_\allowbreak{}lshrti3} & \texttt{u128(u128, int)}\\
\bottomrule
\end{longtable}
\endgroup

\BedrockTableCaption{\texttt{\_\_int128} and Floating-Point Conversion Helpers}
\begingroup\footnotesize
\setlength{\tabcolsep}{2pt}
\begin{longtable}{@{}p{1.55in}p{1.55in}p{2.45in}@{}}
\toprule
\textbf{Conversion} & \textbf{Helper} & \textbf{Signature}\\
\midrule
\endhead
Signed to binary32 & \texttt{\_\allowbreak{}\_\allowbreak{}floattisf} & \texttt{float(i128)}\\
Signed to binary64 & \texttt{\_\allowbreak{}\_\allowbreak{}floattidf} & \texttt{double(i128)}\\
Unsigned to binary32 & \texttt{\_\allowbreak{}\_\allowbreak{}floatuntisf} & \texttt{float(u128)}\\
Unsigned to binary64 & \texttt{\_\allowbreak{}\_\allowbreak{}floatuntidf} & \texttt{double(u128)}\\
Binary32 to signed & \texttt{\_\allowbreak{}\_\allowbreak{}fixsfti} & \texttt{i128(float)}\\
Binary64 to signed & \texttt{\_\allowbreak{}\_\allowbreak{}fixdfti} & \texttt{i128(double)}\\
Binary32 to unsigned & \texttt{\_\allowbreak{}\_\allowbreak{}fixunssfti} & \texttt{u128(float)}\\
Binary64 to unsigned & \texttt{\_\allowbreak{}\_\allowbreak{}fixunsdfti} & \texttt{u128(double)}\\
\bottomrule
\end{longtable}
\endgroup

\BedrockTableCaption{Complex Arithmetic Helpers}
\begingroup\footnotesize
\setlength{\tabcolsep}{2pt}
\begin{longtable}{@{}p{1.55in}p{1.55in}p{2.45in}@{}}
\toprule
\textbf{Operation} & \textbf{Helper} & \textbf{Signature}\\
\midrule
\endhead
Binary32 multiply & \texttt{\_\allowbreak{}\_\allowbreak{}mulsc3} & \texttt{c32(float, float, float, float)}\\
Binary32 divide & \texttt{\_\allowbreak{}\_\allowbreak{}divsc3} & \texttt{c32(float, float, float, float)}\\
Binary64 multiply & \texttt{\_\allowbreak{}\_\allowbreak{}muldc3} & \texttt{c64(double, double, double, double)}\\
Binary64 divide & \texttt{\_\allowbreak{}\_\allowbreak{}divdc3} & \texttt{c64(double, double, double, double)}\\
\bottomrule
\end{longtable}
\endgroup

The four scalar arguments to a complex helper are the real and imaginary parts
of its two operands, in operand order. Bedrock \texttt{long double} has the
binary64 representation and call classification, so conversions involving
\texttt{long double} use the corresponding \texttt{df} helper and
\texttt{long double \_Complex} multiplication or division uses the
corresponding \texttt{dc3} helper.

\subsection{Pointers and External Objects}
A C function pointer is 64 bits wide and its value is a pre-segment near-call
coordinate. Code and data pointers have the
same size and alignment, while their distinct C type categories preserve
separate code and data address spaces.

A common symbol has the ABI alignment of its declared type. An external object
retains the effective alignment and layout of its declared type, including a
supported packing annotation. Natural aggregate alignment is capped at 16
bytes; an explicit extension owns stronger placement requirements.

\subsection{Calls, Relocations, and Relaxation}
The relocation records the linkage path selected for a C call:

\input{abi/c/documents/fragments/call_relocation_quick_reference.tex}

The assembler and linker may relax an instruction sequence, but relaxation
must preserve the C call interface, symbol identity, and observable control
transfer defined by the original relocation. The relaxed sequence remains part
of the same C calling convention.

\subsection{TLS Relationship}
The Bedrock ELF ABI defines TLS models, symbols, relocations, and the
\texttt{GS0}-relative materialization protocols consumed by C code.

\section{C Memory Model}\label{page:c-memory-model}
This section defines how the baseline C memory model maps to Bedrock memory
operations and ordering instructions.

The ISA memory model owns access atomicity, alignment, instruction-order
selectors, and the global sequentially consistent order. This C ABI owns
atomic-object eligibility and the lowering of C operations into those ISA
contracts.

\subsection{Atomic Object Eligibility}
The native lock-free set consists of integer objects and ordinary 64-bit
pointer objects whose size is 1, 2, 4, or 8 bytes and whose address satisfies
the natural alignment for that size. An unaligned atomic object requires the
target constraint diagnostic below. Accordingly, \texttt{\_\_atomic\_always\_lock\_free} returns true
only for those naturally aligned sizes.

The baseline accepts an atomic object only when it belongs to that native
lock-free set. Any other atomic object---including a floating-point, aggregate,
128-bit integer, unsupported-width, or insufficiently aligned
object---is a target constraint violation.

\subsection{Native Atomic Primitives}
Once the object satisfies the width and alignment requirements above, the
compiler uses the following primitive lowering. The ordering sequences are
defined separately below.

(:c-atomic-lowerings:base.atomic_lowerings.LOAD,base.atomic_lowerings.STORE,base.atomic_lowerings.COMPARE_EXCHANGE,base.atomic_lowerings.EXCHANGE:)

\Needspace{2.6in}
\subsection{C Atomic Memory Order Mapping}
(:c-memory-orders:base.memory_orders.RELAXED,base.memory_orders.CONSUME,base.memory_orders.ACQUIRE,base.memory_orders.RELEASE,base.memory_orders.ACQ_REL,base.memory_orders.SEQ_CST:)

For C compare-exchange, the compiler selects the join of the success and
failure orders through the order lattice, independently of their numeric encodings. Consume is
treated as acquire, acquire joined with release is acquire-release, and
sequential consistency dominates every other order. A failure performs a
read-only operation. Release-sequence effects arise from successful writes,
while a failed
\texttt{CMPXCHG.SEQCST} still participates in the global SC order as an SC
load.

The relaxed C thread fence emits a zero-instruction sequence. Consume,
acquire, release, acquire-release, and sequentially consistent thread fences
emit \texttt{AFENCE}.

\subsection{Volatile, Device, and Executable Memory}
Volatile accesses preserve the observable access count and order. C atomic
operations and fences provide the inter-thread ordering defined above.
The C target interface owns the executable-byte publication operation. A
relocation agent completes that operation before transferring control to
patched instructions.

\Needspace{2.0in}
\subsection{C Fetch RMW Lowering}
(:c-atomic-lowerings:base.atomic_lowerings.FETCH_ADD,base.atomic_lowerings.FETCH_SUB,base.atomic_lowerings.FETCH_AND,base.atomic_lowerings.FETCH_OR,base.atomic_lowerings.FETCH_XOR:)

\clearpage
\section{Calling Convention Examples}\label{page:calling-convention-examples}
The examples apply the normative assignment procedure above.

\subsection{Independent Register Classes and Pair Alignment}
\begingroup\small\ttfamily
\begin{tabularx}{\linewidth}{@{}X@{}}
struct pair \{ uint64\_t lo, hi; \};\tabularnewline
uint64\_t mixed(uint64\_t count, double scale,\tabularnewline
\hspace{1.2em}unsigned \_\_int128 wide, struct pair pair, float bias);\tabularnewline
\end{tabularx}
\endgroup
The general cursor begins at 0 and the floating-point cursor begins at 0.
\texttt{count} consumes \texttt{R0}; \texttt{scale} independently consumes
\texttt{F0}. The GENERAL-PAIR argument rounds the general cursor from 1 to 2,
so \texttt{wide} uses \texttt{R3:R2} and \texttt{R1} remains unused. The
address of the caller-owned \texttt{pair} copy uses \texttt{R4}, and
\texttt{bias} uses \texttt{F1}. The stack argument area has size zero.

\subsection{Independent Scalable Register Classes}
Assume \texttt{vec} and \texttt{pred} denote direct scalable vector and
predicate extension types. For
\texttt{vec blend(vec left, pred mask, uint64\_t count, vec right, pred
select, double scale)}, \texttt{left} and \texttt{right} use \texttt{V0} and
\texttt{V1}; \texttt{mask} and \texttt{select} independently use \texttt{P0}
and \texttt{P1}; \texttt{count} uses \texttt{R0}; and \texttt{scale} uses
\texttt{F0}. The direct vector result uses \texttt{V0}. Each class advances
its own cursor exclusively.

\subsection{sret Changes Ordinary Argument Assignment}
\begingroup\small\ttfamily
\begin{tabularx}{\linewidth}{@{}X@{}}
struct triple \{ uint64\_t x, y, z; \};\tabularnewline
struct triple build(uint64\_t tag, signed \_\_int128 wide,\tabularnewline
\hspace{1.2em}double factor);\tabularnewline
\end{tabularx}
\endgroup
Because the result is 24 bytes, its buffer address occupies \texttt{R0} and
the general cursor begins at 1. Thus \texttt{tag} uses \texttt{R1},
\texttt{wide} uses \texttt{R3:R2}, and \texttt{factor} uses \texttt{F0}.
The callee stores the result through the entry value of \texttt{R0} and returns
the same address in \texttt{R0}.

\subsection{Pair Exhaustion of the General Class}
\begingroup\small\ttfamily
\begin{tabularx}{\linewidth}{@{}X@{}}
void pressure(uint64\_t a0, uint64\_t a1, uint64\_t a2,\tabularnewline
\hspace{1.2em}uint64\_t a3, uint64\_t a4, uint64\_t a5, uint64\_t a6,\tabularnewline
\hspace{1.2em}unsigned \_\_int128 wide, uint64\_t tail);\tabularnewline
\end{tabularx}
\endgroup
Arguments \texttt{a0} through \texttt{a6} use \texttt{R0} through
\texttt{R6}. The next even pair would begin past \texttt{R7}, so
\texttt{wide} is placed completely at \texttt{[SP+16]} and the GENERAL class
becomes exhausted. \texttt{tail} therefore uses \texttt{[SP+32]}, while
\texttt{R7} remains unassigned.

\subsection{Variadic Promotions and Slot Traversal}
\begingroup\small\ttfamily
\begin{tabularx}{\linewidth}{@{}X@{}}
int trace(const char *format, ...);\tabularnewline
trace("example", (signed char)-1, (float)1.5, pair);\tabularnewline
\end{tabularx}
\endgroup
The named \texttt{format} argument uses \texttt{R0}. The signed character is
promoted to \texttt{int} and stored at \texttt{[SP+16]}; the \texttt{float} is
promoted to \texttt{double} and stored at \texttt{[SP+32]}. The caller creates
the aggregate copy and stores its address at \texttt{[SP+48]}. A
\texttt{va\_list} initialized by \texttt{va\_start} visits those three slots in
that order by adding 16 after each access.

\Needspace{1.8in}
\subsection{Assembly Conventions}
The assembly examples illustrate the normative rules. They present the
directives, comments, and labels that determine the shown control flow. The
scalar leaf example presents its one-instruction counted loop with \texttt{REP}; the counter supplies
N directly and zero annuls the body. If an implementation sequence conflicts with a
normative rule, the rule controls.

The declarations below state the ABI interfaces used by the examples.

% ABI-ASM-EXAMPLE: scalar-leaf
\subsection{Scalar Leaf Function}
\begingroup\small\ttfamily
\begin{tabularx}{\linewidth}{@{}X@{}}
int sum(int *p, int n);\tabularnewline
\end{tabularx}
\endgroup

On entry, \texttt{p} is in \texttt{R0} and the sign-extended value of
\texttt{n} is in \texttt{R1}. The function uses only volatile
registers, preserves the entry \texttt{SP}, and returns the integer sum in
\texttt{R0}.

\begin{BedrockFormBlock}{2.45in}
\begingroup\small\ttfamily
\begin{tabularx}{\linewidth}{@{}X@{}}
sum:\tabularnewline
\hspace{1.2em}mov.q\hspace{0.9em}r0, r2\tabularnewline
\hspace{1.2em}clr.q\hspace{0.9em}r0\tabularnewline
\hspace{1.2em}maxs.q\hspace{0.65em}r0, r1\tabularnewline
\hspace{1.2em}rep\hspace{1.05em}r1, add.l [r2++], r0\tabularnewline
\hspace{1.2em}ret\tabularnewline
\end{tabularx}
\endgroup
\BedrockFigureCaption{Scalar leaf function}
\end{BedrockFormBlock}

% ABI-ASM-EXAMPLE: nonleaf-preservation
\subsection{Non-Leaf Preservation and Call Alignment}
\begingroup\small\ttfamily
\begin{tabularx}{\linewidth}{@{}X@{}}
int mf\_pipeline(int *tmp, int *dst, int *src, int n,\tabularnewline
\hspace{1.2em}int ca, int cb, int cc);\tabularnewline
\end{tabularx}
\endgroup

The seven arguments occupy \texttt{R0} through \texttt{R6} in source order.
The function keeps \texttt{n}, \texttt{dst}, \texttt{tmp}, and the
first intermediate result live in nonvolatile \texttt{R8} through
\texttt{R11}. It therefore saves and restores both canonical register pairs.

\Needspace{5.0in}
\begingroup\scriptsize\ttfamily
\begin{tabularx}{\linewidth}{@{}X@{}}
mf\_pipeline:\tabularnewline
\hspace{1.2em}pushp\hspace{0.9em}4\tabularnewline
\hspace{1.2em}pushp\hspace{0.9em}5\tabularnewline
\hspace{1.2em}sub.q\hspace{0.9em}8, sp\tabularnewline
\hspace{1.2em}mov.q\hspace{0.9em}r3, r8\tabularnewline
\hspace{1.2em}mov.q\hspace{0.9em}r1, r9\tabularnewline
\hspace{1.2em}mov.q\hspace{0.9em}r0, r10\tabularnewline
\hspace{1.2em}mov.q\hspace{0.9em}r2, r1\tabularnewline
\hspace{1.2em}mov.q\hspace{0.9em}r8, r2\tabularnewline
\hspace{1.2em}mov.q\hspace{0.9em}r4, r3\tabularnewline
\hspace{1.2em}mov.q\hspace{0.9em}r5, r4\tabularnewline
\hspace{1.2em}mov.q\hspace{0.9em}r6, r5\tabularnewline
\hspace{1.2em}call\hspace{1.05em}mf\_fir\_three\tabularnewline
\hspace{1.2em}mov.q\hspace{0.9em}r0, r11\tabularnewline
\hspace{1.2em}mov.q\hspace{0.9em}r9, r0\tabularnewline
\hspace{1.2em}mov.q\hspace{0.9em}r10, r1\tabularnewline
\hspace{1.2em}mov.q\hspace{0.9em}r8, r2\tabularnewline
\hspace{1.2em}call\hspace{1.05em}mf\_scale\_store\tabularnewline
\hspace{1.2em}add.l\hspace{0.9em}r11, r0\tabularnewline
\hspace{1.2em}mov.l\hspace{0.55em}r0, r2\tabularnewline
\hspace{1.2em}mov.q\hspace{0.9em}r9, r0\tabularnewline
\hspace{1.2em}mov.q\hspace{0.9em}r8, r1\tabularnewline
\hspace{1.2em}add.q\hspace{0.9em}8, sp\tabularnewline
\hspace{1.2em}popp\hspace{1.0em}5\tabularnewline
\hspace{1.2em}popp\hspace{1.0em}4\tabularnewline
\hspace{1.2em}jmp\hspace{1.15em}mf\_bias\_sum\tabularnewline
\end{tabularx}
\endgroup
\BedrockFigureCaption{Non-leaf function with a tail transfer}

Each \texttt{PUSHP} saves 16 bytes, so the pair saves leave the entry alignment
unchanged. The additional 8-byte subtraction establishes \texttt{SP mod 16 =
8} before each \texttt{CALL}; the pushed return address therefore gives the
callee a 16-byte-aligned entry stack. The epilogue removes the padding and
restores \texttt{R10:R11} before \texttt{R8:R9}. The final \texttt{JMP} is a
legal tail transfer because the frame and nonvolatile state have already been
restored, the outgoing arguments have the ordinary ABI layout, and the target's
\texttt{R0} result is also the result required by \texttt{mf\_pipeline}.

\end{document}
